\section{Ecosystem context, proxy risks, and case-level audits}\label{si:context}

\subsection{Peru and Chile: a plausible confound is not an identified cause}

The seven stocks assigned \(\Tr=0.5\) years in the ranked test are Peruvian or Chilean anchovy stocks with annual quotas and group-described in-season responsive closures. The code does not verify each stock's specific closure trigger, decision date, or state-assessment cadence. The shorter proxy is therefore a regional group classification. It is not a measured stock-specific treatment.

A retained anchoveta catch series and prior analysis notes associate a roughly 3.7-year feature with ENSO-related environmental variation. That association motivates a confound check, but the v53 ranked test does not re-estimate a causal ENSO model on a fully verified primary-source governance panel. The per-stock AR(1) surrogate reproduces fitted red-noise persistence but not quasi-periodic forcing. Consequently, the observed negative \(\Tr\)-rank association and the higher two-group null-excess in Peru/Chile are compatible with regional forcing, cohort composition, assessment differences, or other omitted factors; they do not establish which factor produced the contrast. The correct result remains: no identifiable approach gradient, plus an opposite-direction proxy-group association with unresolved confounding.

This matters because ENSO and longer regime shifts can structure anchoveta abundance, and oceanographic change has been documented off Peru (Ch\'avez et al., 2003; Guti\'errez et al., 2009). Those references provide ecological context, not an instrumental variable or proof that ENSO caused the v53 contrast.

\subsection{Icelandic cod and haddock: numerical audit of an earlier forensics note}

A recovered cod/haddock audit used current-vintage ICES standard-graph source-data tables, covering cod 1955--2026 and haddock 1979--2026. For 1996--2023, cod SSB has \(n=28\), CV 0.394, catch CV 0.150, fishing-mortality CV 0.199, and recruitment CV 0.208. The previously archived cod CV of 0.387 is confirmed as the same post-1995 concept under an earlier vintage; the current-vintage value is 0.394. For haddock over the same years, SSB CV is 0.359 and recruitment CV is 0.894. The earlier value 0.143 for haddock is not reproduced under standard SSB-CV definitions: detrended-residual CV is 0.345, log-difference volatility is 0.171, and the observed SSB CV is 0.359. That number is therefore not used in v53.

A post-1995 detrended cod spectrum has its strongest peak near 19.58 years; the 10--15-year band contains 14.9\% of the power under the retained calculation. The peak is not a stable, narrow institutional signature and the series is trend-dominated. The numerical audit confirms the earlier coding note's broad environmental/recruitment confound, but neither this CV comparison nor the spectral summary identifies an HCR effect. The haddock recruitment variability result is descriptive and does not adjudicate the governor's review cycle.

\subsection{Northern cod ecosystem evidence}

The northern-cod SSB and mortality values used in Fig. 3 and Table~\ref{tab:cod-values} are copied from the DFO CSAS 2016/026 assessment record. The 1991--1995 crash window has assessment SSB values 735, 382, 101, 31, and 10 kt for 1991--1995, respectively; the full transcribed table is in Supplementary S13.3. The post-1995 portion contains later declines and recoveries, so it is called a non-recovery \emph{record}, not a monotone trend. Tam and Bundy (2019) provide mass-balance descriptions of harp-seal and capelin changes on the Newfoundland and Labrador Shelf; those are ecosystem context only. No seal-predation mortality rate is estimated here, and the constrained-\(M\) decomposition is not reproduced.

\subsection{Model-class diagnostic for scalar autonomous surplus production}

A simple scalar autonomous surplus-production equation with constant removals has the form
\[
\dot S=f(S)-C,\qquad f(S)=\frac{r}{K-\mathfrak s}(K-S)(S-\mathfrak s),
\]
where \(0<\mathfrak s<K\), \(C\geq0\), and \(S\) is the modeled stock. The phase-line shape of \(f\) limits the equilibrium structure. If \(C>0\), the vector field at zero is \(-C\), so the state requires an explicit nonnegativity constraint or a different low-stock removal rule. For extra mortality \(M_x\geq0\), write \(f_M(S)=f(S)-M_xS=S(h(S)-M_x)\) with \(h(S)=f(S)/S\). The zero equilibrium remains; positive equilibria solve \(h(S)=M_x\). When \(0<M_x<\max h\), the quadratic \(h\) can have two positive roots, including a lower threshold shifted above \(\mathfrak s\); at the maximum the roots coalesce and above it positive equilibria disappear.

This is a model-class observation, not an empirical test of northern cod. The assessment series is not an error-free trajectory and removals and mortality are not constant, so no contradiction between this scalar model and the real stock is claimed. The lemma says what this class would imply if its assumptions held; it does not establish those assumptions.

\subsection{Interpretive guardrails}

The contexts above are used to motivate competing mechanisms and the need for discriminating data. They do not turn the spectral screen, the review-interval proxy, a coefficient of variation, or a visual event sequence into evidence of institutional causation. The classification of an environmental mechanism as ``plausible'' is intentionally weaker than a causal finding.

